\documentclass[11pt]{article}
% ===========================================================================
%  Shared preamble for all MPM2D worksheets — input right after \documentclass.
%  (Files beginning with "_" are skipped by mpm2d-worksheet-publish.mjs.)
% ===========================================================================
\usepackage[letterpaper,top=2.2cm,bottom=1.9cm,left=1.6cm,right=1.6cm,headheight=28pt,footskip=24pt]{geometry}
\usepackage{amsmath,amssymb}
\usepackage[T1]{fontenc}
\usepackage[utf8]{inputenc}
\usepackage{lmodern}
\usepackage{xcolor}
\usepackage[most]{tcolorbox}
\usepackage{enumitem}
\usepackage{multicol}
\usepackage{fancyhdr}
\usepackage{tikz}
\usetikzlibrary{calc}
\usepackage{pgfplots}
\pgfplotsset{compat=1.18}

\definecolor{exblue}{HTML}{4A90E2}
\definecolor{exbg}{HTML}{E6F3FF}
\definecolor{qorange}{HTML}{E69138}
\definecolor{qbg}{HTML}{FFF7CC}
\definecolor{keygray}{HTML}{475569}
\definecolor{keybg}{HTML}{F1F5F9}
\definecolor{iagreen}{HTML}{14653B}
\definecolor{titleblue}{HTML}{1D4ED8}
% Rotating palette so each worked example gets its own colour.
\definecolor{ex0}{HTML}{2563A0}\definecolor{ex1}{HTML}{3B7D3B}\definecolor{ex2}{HTML}{8A5A00}
\definecolor{ex3}{HTML}{A3327A}\definecolor{ex4}{HTML}{6D28A3}\definecolor{ex5}{HTML}{B08900}
\definecolor{ex6}{HTML}{0E7490}\definecolor{ex7}{HTML}{9A3412}\definecolor{ex8}{HTML}{15803D}

\pagestyle{fancy}
\fancyhf{}
\fancyhead[L]{\footnotesize NAME: \underline{\hspace{3.4cm}}}
\fancyhead[C]{\textbf{Integration Academy}\\[-2pt]{\scriptsize Math Simplified \textperiodcentered{} Success Amplified}}
\fancyhead[R]{\footnotesize MPM2D}
\fancyfoot[L]{\scriptsize dr.merie@IntegrationAcademy.ca}
\fancyfoot[C]{\scriptsize\thepage}
\fancyfoot[R]{\scriptsize IntegrationAcademy.ca}
\renewcommand{\headrulewidth}{1.2pt}
\renewcommand{\footrulewidth}{0.4pt}
\setlength{\parindent}{0pt}
\setlength{\parskip}{4pt}

% Examples are NOT breakable, so each example + its solution stays on one page.
\newtcolorbox{example}[2]{colframe=#1,colback=#1!7,coltitle=white,colbacktitle=#1,fonttitle=\bfseries,title={#2},arc=3pt,boxrule=1pt}
\newtcolorbox{qbox}[1]{colframe=qorange,colback=qbg,coltitle=white,colbacktitle=qorange,fonttitle=\bfseries,title={#1},arc=3pt,breakable,boxrule=1pt}
\newtcolorbox{keybox}{colframe=keygray,colback=keybg,coltitle=white,colbacktitle=keygray,fonttitle=\bfseries,title={$\checkmark$\ Answer Key},arc=3pt,breakable,boxrule=1pt}
\newtcolorbox{recapbox}{colframe=keygray,colback=keybg,coltitle=white,colbacktitle=keygray,fonttitle=\bfseries,title={Question Recap},arc=3pt,breakable,boxrule=1pt}
\newtcolorbox{ideabox}{colframe=titleblue,colback=titleblue!6,coltitle=white,colbacktitle=titleblue,fonttitle=\bfseries,title={The Big Ideas},arc=3pt,breakable,boxrule=1.2pt}
\newtcolorbox{learnbox}{colframe=iagreen,colback=iagreen!5,coltitle=white,colbacktitle=iagreen,fonttitle=\bfseries,title={Concept Essentials},arc=3pt,breakable,boxrule=1.2pt}
\newcommand{\lh}[1]{\par\smallskip{\bfseries\color{iagreen}#1}\par\nobreak\smallskip}

\newcommand{\soln}{\par\smallskip\textit{\textbf{Solution.}}\ }
\newcommand{\workspace}[1]{\par\vspace{3pt}\fcolorbox{black!30}{white}{\begin{minipage}[t][#1][t]{\dimexpr\linewidth-2\fboxsep\relax}\ \end{minipage}}\par}
\newcommand{\sech}[1]{\par\vspace{8pt}{\Large\bfseries\color{iagreen}#1}\par\nobreak\vspace{2pt}{\color{black!20}\hrule height1pt}\vspace{6pt}}

% A compact coordinate plot sized for an example box: \eplot{xmin}{xmax}{ymin}{ymax}{body}
\newcommand{\eplot}[5]{\begin{center}\begin{tikzpicture}
\begin{axis}[width=6.8cm,height=4.4cm,axis lines=middle,xlabel={\tiny$x$},ylabel={\tiny$y$},
grid=both,grid style={gray!16},xmin=#1,xmax=#2,ymin=#3,ymax=#4,
xtick distance=2,ytick distance=2,tick label style={font=\tiny},axis line style={-},samples=120]
#5
\end{axis}\end{tikzpicture}\end{center}}

% A sine-curve plot (degrees on x, 0..360): \splot{ymin}{ymax}{body}
\newcommand{\splot}[3]{\begin{center}\begin{tikzpicture}
\begin{axis}[width=8.6cm,height=4.6cm,axis lines=middle,xlabel={\tiny$x\,(^\circ)$},ylabel={\tiny$y$},
grid=both,grid style={gray!16},xmin=0,xmax=360,ymin=#1,ymax=#2,
xtick distance=90,ytick distance=1,tick label style={font=\tiny},axis line style={-},samples=180]
#3
\end{axis}\end{tikzpicture}\end{center}}

% A small coordinate plot: \plot{xmin}{xmax}{ymin}{ymax}{body}
\newcommand{\plot}[5]{\begin{center}\begin{tikzpicture}
\begin{axis}[width=8.4cm,height=6cm,axis lines=middle,xlabel={\scriptsize$x$},ylabel={\scriptsize$y$},
grid=both,grid style={gray!16},xmin=#1,xmax=#2,ymin=#3,ymax=#4,
xtick distance=2,ytick distance=2,tick label style={font=\tiny},axis line style={-}]
#5
\end{axis}\end{tikzpicture}\end{center}}

% Right triangle (angle theta at bottom-left, right angle at bottom-right):
%   \rtri{drawBase}{drawHeight}{oppLabel}{adjLabel}{hypLabel}
\newcommand{\rtri}[5]{\begin{center}\begin{tikzpicture}[scale=0.85]
\coordinate (A) at (0,0);\coordinate (B) at (#1,0);\coordinate (C) at (#1,#2);
\draw[thick,exblue] (A)--(B)--(C)--cycle;
\draw[black] ($(B)+(-0.32,0)$)--($(B)+(-0.32,0.32)$)--($(B)+(0,0.32)$);
\node[below] at ($(A)!0.5!(B)$){#4};
\node[right] at ($(B)!0.5!(C)$){#3};
\node[above left] at ($(A)!0.5!(C)$){#5};
\node at (0.7,0.22){$\theta$};
\end{tikzpicture}\end{center}}

% General (oblique) triangle with vertex labels and opposite-side labels:
%   \gtri{Alabel}{Blabel}{Clabel}{aSide}{bSide}{cSide}
\newcommand{\gtri}[6]{\begin{center}\begin{tikzpicture}[scale=0.85]
\coordinate (A) at (0,0);\coordinate (B) at (5,0);\coordinate (C) at (1.6,3);
\draw[thick,exblue] (A)--(B)--(C)--cycle;
\node[below left] at (A){#1};\node[below right] at (B){#2};\node[above] at (C){#3};
\node[below] at ($(A)!0.5!(B)$){#6};
\node[right] at ($(B)!0.5!(C)$){#4};
\node[left] at ($(A)!0.5!(C)$){#5};
\end{tikzpicture}\end{center}}

\fancyhead[R]{\footnotesize MCR3U \textperiodcentered{} 4.3}
\begin{document}

\begin{center}
{\LARGE\bfseries\color{titleblue}Worksheet 4.3 --- Reciprocal Ratios \& Trigonometric Identities}\\[2pt]
{\small Grade 11 Functions, University (MCR3U) \textperiodcentered{} Unit 4: Trigonometric Functions}
\end{center}

\begin{ideabox}
The reciprocal ratios and the Pythagorean and quotient identities let you simplify and prove.
\begin{itemize}[leftmargin=*,itemsep=2pt]
  \item $\csc\theta=\dfrac{1}{\sin\theta}$, $\sec\theta=\dfrac{1}{\cos\theta}$, $\cot\theta=\dfrac{1}{\tan\theta}$.
  \item Quotient: $\tan\theta=\dfrac{\sin\theta}{\cos\theta}$. Pythagorean: $\sin^2\theta+\cos^2\theta=1$.
  \item To prove an identity, rewrite everything in sine and cosine.
\end{itemize}
\end{ideabox}

\begin{learnbox}
\lh{The reciprocal ratios}$\csc\theta=\dfrac{1}{\sin\theta}$, $\sec\theta=\dfrac{1}{\cos\theta}$, and $\cot\theta=\dfrac{1}{\tan\theta}$ complete the six trig ratios.
\lh{Fundamental identities}The identities $\sin^2\theta+\cos^2\theta=1$ and $\tan\theta=\dfrac{\sin\theta}{\cos\theta}$ link the ratios and let you rewrite expressions.
\lh{Proving identities}Work one side into the other using known identities and algebra — a proof is a chain of equal expressions, not an equation to solve.
\end{learnbox}

\sech{Worked Examples}

\begin{example}{ex0}{Example 1 --- Reciprocal}
$\sin\theta=\tfrac35$. Find $\csc\theta$.\soln $\csc\theta=\dfrac{1}{\sin\theta}=\tfrac53$.
\end{example}

\begin{example}{ex1}{Example 2 --- Reciprocal}
$\cos\theta=\tfrac45$. Find $\sec\theta$.\soln $\sec\theta=\dfrac{1}{\cos\theta}=\tfrac54$.
\end{example}

\begin{example}{ex2}{Example 3 --- Reciprocal}
$\tan\theta=2$. Find $\cot\theta$.\soln $\cot\theta=\dfrac{1}{\tan\theta}=\tfrac12$.
\end{example}

\begin{example}{ex3}{Example 4 --- Evaluate}
Find $\sec60^\circ$.\soln $\dfrac{1}{\cos60^\circ}=\dfrac{1}{1/2}=2$.
\end{example}

\begin{example}{ex4}{Example 5 --- Evaluate}
Find $\csc30^\circ$.\soln $\dfrac{1}{\sin30^\circ}=\dfrac{1}{1/2}=2$.
\end{example}

\begin{example}{ex5}{Example 6 --- Evaluate}
Find $\cot45^\circ$.\soln $\dfrac{1}{\tan45^\circ}=\dfrac11=1$.
\end{example}

\begin{example}{ex6}{Example 7 --- Pythagorean}
$\sin\theta=\tfrac35$, acute. Find $\cos\theta$.\soln $\cos^2\theta=1-\tfrac{9}{25}=\tfrac{16}{25}\Rightarrow\cos\theta=\tfrac45$ (a 3-4-5 triangle):\rtri{4}{3}{$3$}{$4$}{$5$}
\end{example}

\begin{example}{ex7}{Example 8 --- Pythagorean}
$\cos\theta=\tfrac{5}{13}$, acute. Find $\sin\theta$.\soln $\sin^2\theta=1-\tfrac{25}{169}=\tfrac{144}{169}\Rightarrow\sin\theta=\tfrac{12}{13}$.
\end{example}

\begin{example}{ex8}{Example 9 --- Prove}
Prove $\tan\theta\cos\theta=\sin\theta$.\soln $\tan\theta\cos\theta=\dfrac{\sin\theta}{\cos\theta}\cdot\cos\theta=\sin\theta$. \checkmark
\end{example}

\clearpage
\sech{Your Turn --- Show Your Work}
For each question, show all steps clearly in the space provided.

\begin{qbox}{Question 1}
$\sin\theta=\tfrac12$: find $\csc\theta$.\workspace{2.4cm}
\end{qbox}

\begin{qbox}{Question 2}
$\cos\theta=\tfrac13$: find $\sec\theta$.\workspace{2.4cm}
\end{qbox}

\begin{qbox}{Question 3}
$\tan\theta=4$: find $\cot\theta$.\workspace{2.4cm}
\end{qbox}

\begin{qbox}{Question 4}
Find $\sec45^\circ$.\workspace{2.4cm}
\end{qbox}

\begin{qbox}{Question 5}
Find $\csc60^\circ$.\workspace{2.4cm}
\end{qbox}

\begin{qbox}{Question 6}
Find $\cot30^\circ$.\workspace{2.4cm}
\end{qbox}

\begin{qbox}{Question 7}
$\sin\theta=\tfrac{8}{17}$, acute: find $\cos\theta$.\workspace{2.4cm}
\end{qbox}

\begin{qbox}{Question 8}
$\cos\theta=\tfrac{24}{25}$, acute: find $\sin\theta$.\workspace{2.4cm}
\end{qbox}

\begin{qbox}{Question 9}
Simplify $\csc\theta\cdot\sin\theta$.\workspace{2.4cm}
\end{qbox}

\begin{qbox}{Question 10}
Simplify $\dfrac{\cos\theta}{\sin\theta}$.\workspace{2.4cm}
\end{qbox}

\begin{qbox}{Question 11}
Evaluate $\sec0^\circ$.\workspace{2.4cm}
\end{qbox}

\begin{qbox}{Question 12}
Prove $\cot\theta\cdot\tan\theta=1$.\workspace{2.6cm}
\end{qbox}

\begin{qbox}{Question 13 --- Challenge}
$\sin\theta=\tfrac23$, acute: find $\tan\theta$ exactly.\workspace{3cm}
\end{qbox}

\clearpage
\sech{Question Recap \& Answer Key}

\begin{recapbox}
\begin{multicols}{2}
\small
\begin{enumerate}[leftmargin=*,itemsep=3pt]
  \item $\sin\theta=\tfrac12$: find $\csc\theta$.
  \item $\cos\theta=\tfrac13$: find $\sec\theta$.
  \item $\tan\theta=4$: find $\cot\theta$.
  \item Find $\sec45^\circ$.
  \item Find $\csc60^\circ$.
  \item Find $\cot30^\circ$.
  \item $\sin\theta=\tfrac{8}{17}$, acute: find $\cos\theta$.
  \item $\cos\theta=\tfrac{24}{25}$, acute: find $\sin\theta$.
  \item Simplify $\csc\theta\cdot\sin\theta$.
  \item Simplify $\dfrac{\cos\theta}{\sin\theta}$.
  \item Evaluate $\sec0^\circ$.
  \item Prove $\cot\theta\cdot\tan\theta=1$.
  \item $\sin\theta=\tfrac23$, acute: find $\tan\theta$ exactly.
\end{enumerate}
\end{multicols}
\end{recapbox}

\begin{keybox}
\begin{multicols}{2}
\small
\begin{enumerate}[leftmargin=*,itemsep=3pt]
  \item $2$
  \item $3$
  \item $\tfrac14$
  \item $\sqrt2$
  \item $\tfrac{2}{\sqrt3}$
  \item $\sqrt3$
  \item $\tfrac{15}{17}$
  \item $\tfrac{7}{25}$
  \item $1$
  \item $\cot\theta$
  \item $1$
  \item true (reciprocals multiply to $1$)
  \item $\cos\theta=\tfrac{\sqrt5}{3}$, $\tan\theta=\tfrac{2}{\sqrt5}$
\end{enumerate}
\end{multicols}
\end{keybox}

\end{document}
